\scanpage{64}
\mspage{6}
% archive foliation: 66
\[
D\longrightarrow o(D)\qquad (\text{``ombre de $D$''})
\]
qui associe à tt disque conforme \unsure{autour de}
$D$ son ombre dans $T_{X',a}$, n'a rien d'injec-
tif. Les \uline{disques ``concentriques''} à $D$
donnent les $\lambda o(D)$ $(0<\lambda\leq1)$, plus
précisément, $D\unsure{\longrightarrow|-}o(D)$ est compatible aux
homothéties de rapport \deleted{$\lambda>0$} $\lambda\in]0,1]$.

Je m'intéresse en fait au cas
où $X'$ est compacte, et munie d'une
\deleted{\unsure{famille}} partie finie de ``pts marqués'', soit
\deleted{$\unsure{I}$}\inserted{$(a_i)_{i\in I}$}, et où on prend une famille
de disques $D_i$ $(i\in I)$ \inserted{\uline{mutuellement disjoints}} autour des $a_i$,
donnant des ombres
\[
o(D_i)\subset T_{X',a_i}.
\]
On cherche des conditions particulières
\inserted{plausibles} sur le système des $D_i$, impliquant
(qui sait ?) qu'il est déterminé par
le système des ombres.
